\section{Methodology}
\label{sec:methodology}

\subsection{Literature and feature evidence}
The study pairs a requirements-based specification of the stable release with source analysis and measurements from the available beta. The stable workflows are drawn from the project requirements, architecture specification, and the specified transfer-password lifecycle. The literature review focused on the functions and design decisions in those workflows. Searches combined terms for tourist trip design, mobile guides, disconnected replicas, content-addressed storage, filesystem persistence, and secure updates. Twenty-one journal or conference papers and one university research report were analyzed. For each work, the review recorded the method or finding and the specific TripyTravy claim it supports or bounds. Author copies and publisher records were used where available. The bibliography includes product manuals and standards as separate supporting sources.

For the feature analysis, the application entry points and the paths they call were traced from authoring through storage, export, import, and presentation. Desktop records were checked against their mapped \texttt{TripPackage} fields to distinguish available inputs from what the mobile reader receives. The Android manifest and Kotlin activity establish the intended platform handoffs. The Dart adapter wired into each application establishes the active container design. The independent Rust code and its tests establish the prototype design. Travefy features were checked against its public builder guide and help articles. This reflects documented Travefy behavior. It provides no account-level test or view of the vendor's implementation.

\subsection{Functional tests and controlled fixtures}
The evaluation ran on Windows 11 with Flutter 3.47.7 and Dart 3.13.5. The tested code was based on commit \texttt{65db5df} with the revision changes included in the accompanying patch. The desktop and mobile Flutter suites ran on the host. Native build artifacts came from the Windows Creator and Android Viewer. The study gathered no build or runtime data for macOS, Linux, or iOS. The Dart core tests used generated trips and local repositories, and the Rust tests exercised the native container in isolation.

A standalone Dart probe used a fresh disk-backed repository to import a legacy Kashmir export. It checked an item on the local checklist, exported a revised trip with a permanent ID, and imported the revision. It asserted that one trip remained, its title reflected the revision, and the checklist item stayed checked. The probe reopened the original file and checked that the newer revision remained current.

For the size test, the probe generated two, four, or eight distinct document records with 1 MiB payloads. Within each pair, the shared case reused exactly the same payload, and the distinct case varied its first byte for each record. Trip and document fields were held constant except for values derived from the bytes. The probe measured the final package length and asserted that the intended record and unique-blob counts survived import. A separate wrapper test packed a small JSON object, mutated one payload byte, and checked whether decoding accepted the changed object.
